\documentclass[11pt]{article}

\usepackage[margin=1.05in]{geometry}
\usepackage[T1]{fontenc}
\usepackage{lmodern}
\usepackage{amsmath,amssymb,amsthm}
\usepackage{booktabs,longtable,array}
\usepackage{listings}
\usepackage{xcolor}
\usepackage{microtype}
\usepackage[hidelinks]{hyperref}
\emergencystretch=1.5em

\newtheorem{theorem}{Theorem}[section]
\newtheorem{proposition}[theorem]{Proposition}
\newtheorem{lemma}[theorem]{Lemma}
\newtheorem{corollary}[theorem]{Corollary}
\theoremstyle{definition}
\newtheorem{definition}[theorem]{Definition}
\theoremstyle{remark}
\newtheorem{remark}[theorem]{Remark}

% Every non-ASCII character that appears inside a listing must be declared
% here (pdflatex + listings); papers/kp-engine/check_listings_unicode.py
% audits the file against this table.
\lstset{
  basicstyle=\ttfamily\small,
  columns=fullflexible,
  keepspaces=true,
  breaklines=true,
  breakatwhitespace=true,
  literate={∑}{{$\sum$}}1
           {∏}{{$\prod$}}1
           {≤}{{$\le$}}1
           {≥}{{$\ge$}}1
           {≠}{{$\ne$}}1
           {‖}{{$\|$}}1
           {ℝ}{{$\mathbb{R}$}}1
           {ℂ}{{$\mathbb{C}$}}1
           {ℕ}{{$\mathbb{N}$}}1
           {ℤ}{{$\mathbb{Z}$}}1
           {→}{{$\to$}}1
           {↔}{{$\leftrightarrow$}}1
           {∀}{{$\forall$}}1
           {∃}{{$\exists$}}1
           {∈}{{$\in$}}1
           {∉}{{$\notin$}}1
           {⊆}{{$\subseteq$}}1
           {∧}{{$\land$}}1
           {¬}{{$\lnot$}}1
           {φ}{{$\varphi$}}1
           {σ}{{$\sigma$}}1
           {Λ}{{$\Lambda$}}1
           {ε}{{$\varepsilon$}}1
           {β}{{$\beta$}}1
           {Δ}{{$\Delta$}}1
           {μ}{{$\mu$}}1
           {⁻}{{$^{-}$}}1
           {¹}{{$^{1}$}}1
           {'}{{'}}1
           {⊤}{{$\top$}}1
           {∅}{{$\emptyset$}}1
           {⟨}{{$\langle$}}1
           {⟩}{{$\rangle$}}1
           {·}{{$\cdot$}}1
           {↑}{{$\uparrow$}}1
           {₀}{{$_0$}}1
           {ᶜ}{{$^{c}$}}1,
}

\newcommand{\lc}[1]{{\normalfont\path{#1}}}
\newcommand{\file}[1]{{\normalfont\path{#1}}}
\newcommand{\lcmath}[1]{\text{\texttt{\detokenize{#1}}}}
\newcommand{\Pol}{\mathcal{P}}
\newcommand{\incomp}{\not\sim}
\newcommand{\Ur}{\varphi}
\newcommand{\Trees}{\mathcal{T}}
\newcommand{\oracle}{\texttt{[propext, Classical.choice, Quot.sound]}}

\title{A Machine-Checked Tree-Graph Inequality and\\
Koteck\'y--Preiss Cluster-Expansion Engine in Lean~4}
\author{Lluis Eriksson}
\date{September 2, 2026}

\begin{document}
\maketitle

\begin{abstract}
We describe a complete, \texttt{sorry}-free formalization in Lean~4 over
Mathlib of the core analytic--combinatorial engine of abstract polymer
cluster expansions: the Penrose tree-graph inequality
$|\Ur(X)| \le \#\{\text{spanning trees of the incompatibility graph of }X\}$,
proved by a partition scheme built on a greedy breadth-first Penrose map with
Boolean-interval fibers; a family of tree-walk (leaf-removal) bounds; the
Koteck\'y--Preiss convergence theorem in both its uniform-majorant form and its
sharp, weight-respecting form
$\sum_n \bigl(\text{pinned cluster weights}\bigr) \le \|z(c)\|\,e^{a(c)}$ under
the bare criterion $\sum_{Y \incomp X}\|z(Y)\|e^{a(Y)} \le a(X)$; the
Mayer--Ursell inversion $\Xi = \exp(\text{cluster sum})$ for every finite polymer
system with absolutely convergent cluster series; and the consequences the
programme consumes --- restriction to sub-volumes and $Z$-ratio bounds,
activity tilting and exponential size tails, a zero-free activity polydisc,
and two lattice-animal counting theorems.  The development is a 35-module,
14{,}441-line layer (\file{YangMills/KP/}) plus its lattice-side consumers;
every theorem stated here depends on exactly Lean's three standard axioms,
\oracle{}, as recorded in the repository's verification ledger and committed
oracle transcripts, and no theorem carries a project axiom.  We state each
result in ordinary notation next to its Lean declaration, explain the design
decisions (finite polymer types, coordinate-pinned clusters, admissible
parent/level structures as the tree index, exponentials of \texttt{tsum}s
regrouped through a sigma type, volume-uniformity through per-polymer
bounds), tabulate the axiom oracle, and list what is \emph{not} formalized.
To our knowledge --- after a literature and library search --- this is the
first machine-checked proof of the Penrose tree-graph inequality and of the
Koteck\'y--Preiss criterion in any proof assistant; the companion papers on
the volume-uniform area law and on the thermodynamic limit are its first
consumers.
\end{abstract}

\section{Introduction}\label{sec:intro}

Cluster expansions are the workhorse of rigorous statistical mechanics in
perturbative regimes.  Given a family of \emph{polymers} with a hard-core
incompatibility relation and small complex \emph{activities}, the polymer
partition function $\Xi$ is a finite sum over compatible families, and the
cluster expansion expresses $\log \Xi$ as an absolutely convergent series
indexed by \emph{clusters}: tuples of polymers whose incompatibility graph is
connected, weighted by the Ursell (Mayer) coefficient $\Ur$.  Three results
carry the whole theory.  The \emph{tree-graph inequality} of
Penrose~\cite{Penrose1967} bounds $|\Ur|$ by a count of spanning trees; the
\emph{Koteck\'y--Preiss criterion}~\cite{KoteckyPreiss1986} turns that
combinatorial bound into absolute convergence with explicit, per-polymer
constants; and the \emph{Mayer--Ursell inversion} identifies the sum of the
series with $\log \Xi$.  Modern expositions --- Brydges' lectures~\cite{Brydges1986},
Ueltschi's review~\cite{Ueltschi2004}, the sharpened criteria of
Fern\'andez--Procacci~\cite{FernandezProcacci2007} and
Bissacot--Fern\'andez--Procacci~\cite{BFP2010}, Faris' combinatorial
account~\cite{Faris2010}, and Chapter~5 of Friedli--Velenik's
textbook~\cite{FriedliVelenik2017} --- present the material as a small number
of clean theorems.  The proofs, however, are exactly the kind of mathematics
in which errors hide in bookkeeping rather than in ideas: sign-reversing
involutions on edge sets, symmetry factors between ordered tuples and
multisets, factorial weights cancelling against tree counts, and the
interchange of an exponential with an infinite sum.  That makes them a natural
target for machine verification, and also a genuine test of whether current
proof assistants can absorb a piece of constructive statistical mechanics
end to end.

This paper reports such a verification.  Within the author's Lean~4
programme on lattice Yang--Mills theory (repository
\texttt{THE-ERIKSSON-PROGRAMME}; see~\cite{Eriksson2026Master,Eriksson2026Audit}),
the directory \file{YangMills/KP/} contains a 35-module, 14{,}441-line
formalization of the abstract polymer cluster expansion, built over
Mathlib~\cite{mathlib} on Lean~4~\cite{Lean4}.  Two earlier papers in the
programme are \emph{consumers} of this layer: the volume-uniform Wilson-loop
area law~\cite{Eriksson2026AreaLaw} (viXra 2607.0005) and the thermodynamic
limit of local Gibbs states~\cite{Eriksson2026ThermoLimit} (viXra 2607.0091).
Both describe the engine only in passing.  The present paper is the missing
account of the engine itself: precise statements, the design decisions that
made the formalization tractable, and the axiom oracle.

\paragraph{Contributions.}
\begin{enumerate}
\item \textbf{The Penrose tree-graph inequality, machine-checked}
  (\lc{abs_ursell_le_card_spanningTrees}, Section~\ref{sec:penrose}): for every
  polymer system and every tuple $X$, $|\Ur(X)| \le \#\Trees(G_X)$.  The proof
  formalizes the partition-scheme mechanism (Boolean-interval fibers telescope)
  and the concrete greedy breadth-first Penrose scheme, including the fiber
  characterization that is usually left to the reader.
\item \textbf{Tree-walk bounds} (Section~\ref{sec:walk}): an abstract
  leaf-removal lemma \lc{tree_walk_bound} and its vertexwise and root-pinned
  variants, which convert per-tree sums of pair kernels into products of
  one-step budgets.
\item \textbf{Koteck\'y--Preiss convergence in two strengths}
  (Sections~\ref{sec:walk}--\ref{sec:sharp}): the uniform-majorant bound
  $\mathrm{clusterWeight}_n \le (\sum\|z\|)(eA)^n$, and the sharp, per-polymer
  bound $\sum_{n \le N}\mathrm{pinnedClusterWeight}(c,n) \le \|z(c)\|e^{a(c)}$
  under the bare criterion (\lc{kp_pinned_cluster_bound}).  The latter is the
  form that survives the infinite-volume limit.
\item \textbf{The Mayer--Ursell inversion} (\lc{partition_eq_exp_clusterSum},
  Section~\ref{sec:mayer}): $\Xi = \exp(\text{clusterSum})$ for every finite
  polymer system whose cluster series is absolutely convergent, hence under the
  KP criterion; and, as a corollary, non-vanishing of $\Xi$ on a closed activity
  polydisc (\lc{partition_withActivity_ne_zero_of_kp}).
\item \textbf{The volume-uniform toolkit} (Section~\ref{sec:restrict}):
  restriction of a polymer system to a sub-volume with the $Z$-ratio exponent
  bound, activity tilting with exponential size tails, and two lattice-animal
  counting theorems (\lc{card_connectedPolymers_le} for plaquette polymers,
  \lc{animal_card_le} for bounded-degree graphs) that make the criterion's
  smallness window depend only on the dimension.
\end{enumerate}

\paragraph{What is claimed and what is not.}
Everything stated as a theorem below is a Lean declaration whose axiom
oracle prints \oracle{} in the repository's records
(Section~\ref{sec:oracle}); no theorem carries a \texttt{sorry} or a
project axiom.  The mathematics is classical; the contribution is the
mechanization and the explicitness.  Section~\ref{sec:not} lists what is not
formalized: infinite polymer systems, the Fern\'andez--Procacci
sharpening, the Brydges--Kennedy / Abdesselam--Rivasseau tree identities,
the exact Cayley count, and the correlation-function cluster expansion.

\section{Related work}\label{sec:related}

\paragraph{Mathematics.}
Penrose's partition scheme~\cite{Penrose1967} (see also~\cite{Penrose1963})
is the source of the tree-graph inequality in the form used here; the
Boolean-interval formulation follows the modern presentations
in~\cite{FernandezProcacci2007,Faris2010,Scott2005}.  Koteck\'y and
Preiss~\cite{KoteckyPreiss1986} proved convergence by induction on cluster
size under the criterion $\sum_{Y\incomp X}\|z(Y)\|e^{a(Y)} \le a(X)$;
Dobrushin~\cite{Dobrushin1996} and Miracle-Sol\'e~\cite{MiracleSole2000} gave
alternative proofs, Ueltschi~\cite{Ueltschi2004} an efficient one, and
Fern\'andez--Procacci~\cite{FernandezProcacci2007} and
Bissacot--Fern\'andez--Procacci~\cite{BFP2010} sharper criteria.  Brydges'
lectures~\cite{Brydges1986} and Friedli--Velenik~\cite{FriedliVelenik2017}
(Theorem~5.4, Proposition~5.3) are the expository references whose statements
the formalization tracks most closely.  The connection between the polymer
gas and the Lov\'asz local lemma~\cite{ScottSokal2005} is relevant to the
formalization landscape below.

\paragraph{Formalizations.}
We searched (September 2026) the Lean Mathlib library, the Isabelle
Archive of Formal Proofs, the Coq/Rocq opam ecosystem, and the general
literature for a formalization of the tree-graph inequality, the
Koteck\'y--Preiss theorem, the Mayer expansion, or an abstract polymer model.
We found none.  The repository's own planning document
(\file{docs/kp-cluster-expansion-plan.md}, May 2026) already recorded that
Mathlib contains no polymer model, no Ursell coefficient and no cluster
expansion; that remains the case at the pinned Mathlib revision.  The
closest neighbours are formalizations of the probabilistic method and of the
Lov\'asz local lemma (in Isabelle/HOL, Edmonds~\cite{Edmonds2023}), which
share the combinatorial flavour of the independent-set polynomial but do not
touch Ursell coefficients, tree-graph bounds or convergence criteria.  We
therefore believe --- with the usual caveat that absence of evidence in a
search is not proof of absence --- that the results of
Sections~\ref{sec:penrose}--\ref{sec:mayer} are the first machine-checked
proofs of these theorems in any proof assistant.  Within the programme
itself, the Catalan/rooted-tree majorant material of
\cite{Eriksson2026Catalan,Eriksson2026C1} is a separate, already documented
lane (\file{YangMills/KP/RootedCatalan*.lean},
\file{PlaneTreeCatalan.lean}, \file{LabeledPlaneTreeCatalan.lean}); it is
not re-described here.

\section{The abstract polymer system in Lean}\label{sec:setting}

All definitions live in \texttt{namespace YangMills.KP}.  We write $\Pol$ for
the polymer type, $X \incomp Y$ for incompatibility and $z$ for the activity.

\begin{definition}[polymer system; \file{YangMills/KP/Basic.lean}]
A \emph{polymer system} is a structure
\begin{lstlisting}
structure PolymerSystem where
  Polymer : Type*
  incomp : Polymer → Polymer → Prop
  incomp_symm : ∀ X Y, incomp X Y → incomp Y X
  incomp_self : ∀ X, incomp X X
  activity : Polymer → ℂ
\end{lstlisting}
A finite family $S$ is \emph{admissible} (\lc{Admissible}) if no two distinct
members are incompatible, and the \emph{partition function} of a finite region
$\Lambda \subseteq \Pol$ is
\[
  \Xi(\Lambda) \;=\; \sum_{\substack{S \subseteq \Lambda\\ S \text{ admissible}}}
  \;\prod_{X \in S} z(X)
  \qquad (\lcmath{partition}).
\]
\end{definition}

The hard core is built in as reflexivity of $\incomp$; this is the
Koteck\'y--Preiss convention $\gamma \incomp \gamma$ and it is what makes the
criterion sum below contain the $X$-term.  The first nontrivial fact is the
factorization $\Xi(\Lambda_1 \cup \Lambda_2) = \Xi(\Lambda_1)\Xi(\Lambda_2)$ over
disjoint, mutually compatible regions (\lc{partition_union}), proved by an
explicit bijection $S \mapsto (S\cap\Lambda_1, S\cap\Lambda_2)$.

\begin{definition}[clusters and the Ursell coefficient;
\file{Cluster.lean}, \file{Ursell.lean}]
For an $n$-tuple $X : \mathrm{Fin}\,n \to \Pol$ the \emph{incompatibility graph}
$G_X$ (\lc{incompGraph}) is the simple graph on $\mathrm{Fin}\,n$ with
$i \sim j$ iff $i \ne j$ and $X_i \incomp X_j$; $X$ is a \emph{cluster}
(\lc{IsCluster}) iff $G_X$ is connected.  The \emph{Ursell coefficient} is the
signed count of connected spanning subgraphs,
\[
  \Ur(X) \;=\; \sum_{\substack{E \subseteq E(G_X)\\ (\mathrm{Fin}\,n,\,E)\ \text{connected}}}
  (-1)^{|E|} \;\in\; \mathbb{Z}
  \qquad (\lcmath{ursell}).
\]
\end{definition}

\begin{lstlisting}
noncomputable def ursell {n : ℕ} (X : Fin n → P.Polymer) : ℤ := by
  classical
  exact ∑ E ∈ (incompGraph P X).edgeFinset.powerset.filter
      (fun E : Finset (Sym2 (Fin n)) =>
        (SimpleGraph.fromEdgeSet (↑E : Set (Sym2 (Fin n)))).Connected),
    (-1 : ℤ) ^ E.card
\end{lstlisting}

Spanning subgraphs are represented as \texttt{Finset}s of edges
$E \subseteq E(G_X)$, and the graph they span is Mathlib's
\lc{SimpleGraph.fromEdgeSet}.  This choice --- edge finsets on a fixed vertex
type $\mathrm{Fin}\,n$ rather than Mathlib's \texttt{Subgraph} --- is used
consistently through the Penrose scheme and the Mayer inversion, and it is the
single most consequential representation decision in the layer: Boolean
intervals $\{E : T \subseteq E \subseteq R\}$, sign-reversing edge toggles, and
component partitions are all \texttt{Finset} operations.  Sanity checks
proved on the definition: $\Ur = 1$ on singletons, $-1$ on an incompatible
pair, $+2$ on a triangle (\lc{ursell_fin_one}, \lc{ursell_fin_two},
\lc{ursell_fin_three}, the last by kernel \texttt{decide} over the eight
spanning subgraphs), and $\Ur(X) = 0$ whenever $X$ is not a cluster
(\lc{ursell_eq_zero_of_not_isCluster}).  On the complete graph the recurrence
$d(n+1) = -n\,d(n)$ and the closed form $\Ur(K_{n+1}) = (-1)^n n!$ are
theorems (\lc{ursellComplete_recurrence}, \lc{ursellComplete_eq},
\file{UrsellRecurrence.lean}), proved by fibering the vanishing all-subgraphs
sum by the component of vertex~$0$ and killing the large fibers with a
sign-reversing edge-toggle involution.

\begin{definition}[cluster sum, weights, criterion;
\file{Expansion.lean}, \file{ClusterWeight.lean}, \file{PinnedCluster.lean},
\file{Criterion.lean}]
For a \emph{finite} polymer type ($[\texttt{Fintype P.Polymer}]$):
\begin{align}
  \mathrm{clusterSum}(P) &= \sum_{n=0}^{\infty} \frac{1}{(n+1)!}
    \sum_{X : \mathrm{Fin}(n+1) \to \Pol} \Ur(X) \prod_{i} z(X_i)
    \;\in\;\mathbb{C},\label{eq:clustersum}\\
  \mathrm{clusterWeight}(P,n) &= \frac{1}{(n+1)!}
    \sum_{X : \mathrm{Fin}(n+1) \to \Pol} |\Ur(X)| \prod_i \|z(X_i)\|
    ,\\
  \mathrm{pinnedClusterWeight}(P,c,n) &= \frac{1}{(n+1)!}
    \sum_{\substack{X : \mathrm{Fin}(n+1) \to \Pol\\ X_0 = c}}
    |\Ur(X)| \prod_i \|z(X_i)\| .
\end{align}
(in Lean: \lc{clusterSum}, \lc{clusterWeight}, \lc{pinnedClusterWeight}).  A weight $a : \Pol \to \mathbb{R}$ satisfies the \emph{Koteck\'y--Preiss
criterion} (\lc{KPCriterion}) if $a \ge 0$ and
\begin{equation}\label{eq:kp}
  \sum_{Y \,:\, X \incomp Y} \|z(Y)\| \, e^{a(Y)} \;\le\; a(X)
  \qquad \text{for every } X \in \Pol .
\end{equation}
\end{definition}

Three design decisions are visible here.  (i)~\emph{Finite polymer types.}
The outer sum in~\eqref{eq:clustersum} is a genuine \texttt{tsum} over the
cluster size, but every inner sum is a \texttt{Finset} sum; there is no
infinite polymer set anywhere in the engine.  Volume-uniformity is obtained
not by working in infinite volume but by proving bounds whose shape does not
mention $\#\Pol$ (Sections~\ref{sec:sharp} and~\ref{sec:restrict}).
(ii)~\emph{Ordered tuples, not multisets.}  Clusters are functions
$\mathrm{Fin}(n+1) \to \Pol$ with the symmetry factor $1/(n+1)!$; this avoids
multiset combinatorics entirely and makes the Ursell coefficient
relabeling-invariant by a transported graph isomorphism
(\lc{ursell_comp_equiv}).  (iii)~\emph{Coordinate pinning.}  A pinned cluster
is one with $X_0 = c$, not one containing $c$; pinning is then an exact
fiberwise decomposition, $\mathrm{clusterWeight}(P,n) = \sum_c
\mathrm{pinnedClusterWeight}(P,c,n)$ (\lc{clusterWeight_eq_sum_pinned}), with
no overcounting.  (Paper-level remark, not formalized: by symmetry the
coordinate-pinned sum equals the sum over tuples \emph{containing} $c$ weighted
by the fraction of coordinates equal to $c$, so the two conventions differ
only by that fraction.)  Two elementary consequences of the criterion are
recorded once and reused: $\|z(X)\| \le a(X)$ (\lc{kp_activity_le}) and the
unweighted neighbour bound $\sum_{Y\incomp X}\|z(Y)\| \le a(X)$
(\lc{kp_neighbor_sum_le}), both because the $X$-term is in the sum and
$e^{a} \ge 1$.

\section{The Penrose tree-graph inequality}\label{sec:penrose}

\subsection{Statement}

Let $\Trees(G)$ (\lc{spanningTrees}, \file{PenrosePrep.lean}) be the finset of
edge subsets $E \subseteq E(G)$ such that $\mathrm{fromEdgeSet}\,E$ is a tree
on the full vertex set $\mathrm{Fin}\,n$.

\begin{theorem}[Penrose tree-graph inequality]\label{thm:penrose}
\textup{[\lc{abs_ursell_le_card_spanningTrees}, \file{KP/PenroseFiber.lean}]}
For every polymer system $P$, every $n$ and every $X : \mathrm{Fin}\,n \to \Pol$,
\[
  |\Ur(X)| \;\le\; \#\Trees(G_X).
\]
\end{theorem}
\begin{lstlisting}
theorem abs_ursell_le_card_spanningTrees (P : PolymerSystem) {n : ℕ}
    (X : Fin n → P.Polymer) [Fintype (incompGraph P X).edgeSet] :
    |ursell P X| ≤ ((spanningTrees (incompGraph P X)).card : ℤ)
\end{lstlisting}

Two uniform corollaries follow (same file).  With $\mathrm{treeCount}(n)$ the
number of diagonal-free edge sets that are trees on $\mathrm{Fin}\,n$
(\lc{treeCount}), $|\Ur(X)| \le \mathrm{treeCount}(n)$
(\lc{abs_ursell_le_treeCount}), and
$\mathrm{treeCount}(m+1) \le (m+1)^{m+1}$ (\lc{treeCount_le_pow}), hence
$|\Ur(X)| \le (m+1)^{m+1}$ for every $(m+1)$-tuple
(\lc{abs_ursell_le_succ_pow}).  The tree count is \emph{not} Cayley's
$(m+1)^{m-1}$: the injection used is the greedy-parent map
$T \mapsto \mathrm{bfsParent}\,T$ into $\mathrm{Fin}(m+1) \to \mathrm{Fin}(m+1)$,
which is trivial to verify once one knows that a spanning tree is recovered
from its parent function (\lc{penroseTree_of_spanningTree}).  The lost
factor $(m+1)^2$ is absorbed by the analytic inequality
$(n+1)^n \le e^n\, n!$ (\lc{succ_pow_le_exp_mul_factorial},
\file{PenrosePrep.lean}, an induction from $(1+1/k)^k \le e$), which is all
the convergence argument needs.

\subsection{Proof architecture: partition schemes}

The proof is Penrose's, in the Boolean-interval formulation.

\begin{lemma}[interval signed sum]
\textup{[\lc{interval_signed_sum}, \file{PenroseScheme.lean}]}
For finsets $T \subseteq R$,
$\sum_{T \subseteq E \subseteq R} (-1)^{|E|} = (-1)^{|T|}\,[R = T]$.
\end{lemma}

\begin{lemma}[abstract partition scheme]
\textup{[\lc{abs_signedSum_le_of_scheme}]}
Let $S$ be a finset of finsets, $\mathrm{Tr}$ a finset (the ``trees''), and
$\pi, R$ maps of finsets with $\pi(S) \subseteq \mathrm{Tr}$ and, for every
$T \in \mathrm{Tr}$ and every $E$,
$(E \in S \wedge \pi E = T) \iff (T \subseteq E \subseteq R\,T)$.
Then $|\sum_{E \in S} (-1)^{|E|}| \le \#\mathrm{Tr}$.
\end{lemma}

The graph-level statement \lc{abs_ursell_le_card_spanningTrees_of_scheme}
instantiates $S$ to the connected spanning edge subsets of $G_X$ and
$\mathrm{Tr}$ to $\Trees(G_X)$, carrying the scheme $(\pi,R)$ and its two
properties \texttt{hmaps}/\texttt{hfiber} as explicit hypotheses.  This
separation --- an abstract telescoping lemma plus a concrete scheme --- is
what made the development incremental: the abstract half was proved and
oracle-checked weeks before the scheme existed, and the scheme's two
obligations were closed one at a time.

\subsection{The greedy BFS scheme and its fibers}

The concrete scheme (\file{PenroseBFS.lean}, \file{PenroseFiber.lean}) is the
classical greedy breadth-first construction on the vertex set
$\mathrm{Fin}(m+1)$ rooted at $0$:
\begin{itemize}
\item $\mathrm{bfsLevel}\,E\,v$ is the graph distance from $0$ in
  $\mathrm{fromEdgeSet}\,E$ (Mathlib's \lc{SimpleGraph.dist});
\item $\mathrm{bfsParent}\,E\,v$ is the \emph{least} neighbour of $v$ one
  level closer to the root (\lc{Finset.min'} over the candidate set);
\item $\pi(E) = \mathrm{penroseTree}\,E = \{\,\{\mathrm{bfsParent}\,E\,v, v\} : v \ne 0\,\}$;
\item $R(T) = \mathrm{penroseEnvelope}\,H\,T$ is the set of edges $\{u,v\}$ of
  the ambient graph $H$ whose endpoints lie on the same $T$-level, or on
  adjacent levels with the lower endpoint not smaller than the greedy parent
  of the upper one (\lc{penroseRel}).
\end{itemize}
Property \texttt{hmaps} is \lc{penroseTree_mem_spanningTrees}: for connected
$E$ supported on $H$, $\pi(E)$ is connected (induction on the level), has
exactly $m$ edges (the child is the higher endpoint, so $v \mapsto
\{\mathrm{parent}\,v, v\}$ is injective), hence is a tree by Mathlib's
\lc{isTree_iff_connected_and_card}.  Property \texttt{hfiber} is the theorem
usually stated without proof:

\begin{theorem}[Penrose fiber characterization]
\textup{[\lc{penrose_hfiber}]}
For a spanning tree $T$ of $H$ and any edge finset $E$: $E$ is a connected
spanning subgraph of $H$ with $\pi(E) = T$ if and only if
$T \subseteq E \subseteq R(T)$.
\end{theorem}

The proof needs four facts about trees and BFS levels that Mathlib did not
supply directly: distances only shrink in a supergraph
(\lc{dist_le_dist_of_le}); a function that changes by at most one across
edges grows at most by the length along any walk (\lc{level_le_of_walk});
in a tree every path realizes the distance
(\lc{isTree_path_length_eq_dist}, via path uniqueness and the geodesic's
bypass); and adjacent tree vertices lie on consecutive levels
(\lc{isTree_adj_level}).  From these: on the interval $T \subseteq E
\subseteq R(T)$ the BFS levels of $E$ equal those of $T$
(\lc{bfsLevel_eq_of_interval}), the greedy parents agree
(\lc{bfsParent_eq_of_interval}), and $\pi$ fixes spanning trees
(\lc{penroseTree_of_spanningTree}, by parent uniqueness
\lc{isTree_parent_unique}); conversely $\pi$ preserves levels and parents
(\lc{bfsLevel_penroseTree_eq}, \lc{bfsParent_penroseTree_eq}), which places
every edge of a connected $E$ in $R(\pi E)$.  The $n = 0$ case of
Theorem~\ref{thm:penrose} is handled separately ($\Ur = 0$ because no graph
on an empty vertex type is connected in Mathlib's sense).

\section{Tree-walk bounds and the uniform Koteck\'y--Preiss bound}\label{sec:walk}

The bridge from tree counting to activity sums is the leaf-removal induction.
It is stated once, abstractly, over an arbitrary finite vertex type so that
removing a leaf is a plain induction on $n$ (function spaces split at a
coordinate via \lc{Equiv.funSplitAt}) rather than a $\mathrm{Fin}$-reindexing.

\begin{theorem}[tree-walk bound]\label{thm:walk}
\textup{[\lc{tree_walk_bound}, \file{KP/WalkBound.lean}]}
Let $V$ be a finite type with $\#V = n+1$, $r \in V$, $p : V \to V$ a parent
map and $\mathrm{lev} : V \to \mathbb{N}$ with $\mathrm{lev}(p\,v) < \mathrm{lev}(v)$
for $v \ne r$.  Let $\beta$ be a finite type, $I : \beta \to \beta \to
\mathbb{R}_{\ge 0}$, $w : \beta \to \mathbb{R}_{\ge 0}$ and $A \ge 0$ with
$\sum_y I(x,y)\,w(y) \le A$ for all $x$.  Then
\[
  \sum_{X : V \to \beta} \Bigl(\prod_{v \ne r} I(X_{p v}, X_v)\, w(X_v)\Bigr)\, w(X_r)
  \;\le\; A^{n} \sum_{y} w(y).
\]
\end{theorem}
\begin{lstlisting}
theorem tree_walk_bound (n : ℕ) :
    ∀ (V : Type) [Fintype V] [DecidableEq V] (β : Type*) [Fintype β]
      (r : V) (p : V → V) (lev : V → ℕ)
      (I : β → β → ℝ) (w : β → ℝ) (A : ℝ),
      Fintype.card V = n + 1 →
      (∀ v, v ≠ r → lev (p v) < lev v) →
      (∀ x y, 0 ≤ I x y) → (∀ y, 0 ≤ w y) → 0 ≤ A →
      (∀ x : β, ∑ y, I x y * w y ≤ A) →
      ∑ X : V → β,
        (∏ v ∈ Finset.univ.filter (fun v => v ≠ r),
          I (X (p v)) (X v) * w (X v)) * w (X r)
        ≤ A ^ n * ∑ y, w y
\end{lstlisting}

Two variants are proved in the same style.  \lc{tree_walk_bound_vertexwise}
lets the kernel $I_v$ and the budget $A_v$ depend on the child vertex,
concluding $\le (\prod_{v\ne r} A_v)\sum_y w_{\mathrm{root}}(y)$; it is the
form consumed by degree-sensitive leaf summations in the programme's
renormalization-group lane (\file{YangMills/RG/AppendixFSecondUrsellLeafSummation.lean}).
\lc{tree_walk_bound_pinned} (\file{PinnedWalk.lean}) fixes the root value
$X_r = x_0$ and concludes $\le A^n$ with \emph{no} extensive factor
$\sum_y w(y)$.  Its proof is a small trick worth recording: rather than
re-running the induction, extend $\beta$ by a marker \texttt{none} of weight
$T$ that interacts like $x_0$ and can never sit at a non-root vertex; the free
bound gives $T\cdot(\text{pinned sum}) \le A^n (T + \sum w)$ for every
$T \ge 0$, and one large $T$ forces the claim without any limit.

\begin{theorem}[uniform-majorant KP bound]
\textup{[\lc{kp_per_size_bound}, \lc{kp_convergence}, \lc{kp_norm_clusterSum_le}, \file{KPBound.lean}]}
If $a$ satisfies the criterion~\eqref{eq:kp} and $a \le A$ pointwise, then
for every $n$
\[
  \mathrm{clusterWeight}(P,n) \;\le\; \Bigl(\sum_{y}\|z(y)\|\Bigr)\,(eA)^n ,
\]
so for $eA < 1$ the cluster series converges absolutely and
$\|\mathrm{clusterSum}(P)\| \le (\sum_y \|z(y)\|)/(1-eA)$.
\end{theorem}

The proof composes Theorem~\ref{thm:penrose} (expanding $|\Ur(X)|$ over trees
$T$ and dominating $\#\{T \subseteq E(G_X)\}$ by an indicator product over the
edges of $T$), the parent-indexed form of a tree product
(\lc{prod_tree_eq_prod_parents}: $\prod_{e \in T} f(e) = \prod_{v \ne 0}
f(\{\mathrm{parent}\,v, v\})$), the per-tree assignment bound
\lc{tree_assignment_sum_le} obtained from Theorem~\ref{thm:walk} with
$I(x,y) = [x \incomp y]$, $w = \|z\|$ and the budget from
\lc{kp_neighbor_sum_le}, and finally $\mathrm{treeCount}(n+1)/(n+1)! \le e^n$.
This is the Koteck\'y--Preiss theorem in the form ``smallness $eA<1$'', with
$A$ a uniform majorant of the weight.  For a lattice gas with $a(c) = t|c|$ the
only uniform majorant is $A = t\cdot\#\Pol$, which is volume-dependent; the
sharp form of the next section removes $A$ altogether.  A pinned version with
the same defect, $\mathrm{pinnedClusterWeight}(P,c,n) \le \|z(c)\|(eA)^n$
(\lc{kp_pinned_per_size_bound}, \file{PinnedBound.lean}), served as the
stepping stone.

\section{The sharp Koteck\'y--Preiss bound}\label{sec:sharp}

\begin{theorem}[sharp pinned bound]\label{thm:sharp}
\textup{[\lc{kp_pinned_cluster_bound}, \file{KP/SharpShell.lean}]}
If $a$ satisfies the criterion~\eqref{eq:kp}, then for every polymer $c$ and
every truncation $N$,
\[
  \sum_{n=0}^{N} \mathrm{pinnedClusterWeight}(P,c,n) \;\le\; \|z(c)\|\,e^{a(c)} .
\]
Consequently (\lc{pinned_cluster_summable_sharp}) the pinned series is
summable, with sum at most $\|z(c)\|e^{a(c)}$;
the total weights are summable (\lc{kp_clusterWeight_summable_sharp});
the cluster series converges absolutely (\lc{kp_convergence_sharp},
\file{SharpKP.lean}); and
$\|\mathrm{clusterSum}(P)\| \le \sum_c \|z(c)\|e^{a(c)}$
(\lc{kp_norm_clusterSum_le_sharp}).
\end{theorem}
\begin{lstlisting}
theorem kp_pinned_cluster_bound (P : PolymerSystem) [Fintype P.Polymer]
    {a : P.Polymer → ℝ} (h : KPCriterion P a) (c : P.Polymer) (N : ℕ) :
    ∑ n ∈ Finset.range (N + 1), pinnedClusterWeight P c n
    ≤ ‖P.activity c‖ * Real.exp (a c)
\end{lstlisting}

No uniform majorant, no geometric smallness hypothesis, and no quantity of
the shape $\#\Pol$ appears: the right-hand side is a per-polymer quantity.
This is the statement that Friedli--Velenik's Theorem~5.4 and Ueltschi's
review deliver; the formalized route is neither the original
Koteck\'y--Preiss induction on cluster size nor Ueltschi's, but a
\emph{shell decomposition} of rooted trees, chosen (in
\file{docs/SHARP-KP-PLAN.md}) because it reuses Theorem~\ref{thm:penrose} only
through rooted-tree data and keeps every sum finite.

\paragraph{The analytic half (\texttt{SharpMajorant.lean}).}
Define the depth-$D$ majorant by the recursion
\[
  \Phi_0(c) = 1, \qquad
  \Phi_{D+1}(c) = \exp\Bigl(\sum_{c' \,:\, c \incomp c'} \|z(c')\|\,\Phi_D(c')\Bigr)
  \qquad (\lcmath{kpMajorant}).
\]
Under the criterion, $\Phi_D(c) \le e^{a(c)}$ for every depth
(\lc{kpMajorant_le_exp}): induction on $D$, the criterion absorbing one shell
per step.  This is the only place the criterion is used.

\paragraph{The combinatorial half (\texttt{SharpShell.lean}, 3{,}079 lines).}
An \emph{admissible structure} on $\mathrm{Fin}(n+1)$ is a pair $(p,\mathrm{lev})$
of a parent map and a level map into $\mathrm{Fin}(D+1)$ with
$\mathrm{lev}(0) = 0$ and $\mathrm{lev}(p\,v) + 1 = \mathrm{lev}(v)$ for $v \ne 0$
(\lc{IsAdmissible}).  The \emph{raw rooted tree-sum} pinned at $c$ is
\[
  \mathrm{treeSumRaw}(c,D,n) \;=\; \sum_{(p,\mathrm{lev})\ \text{admissible},\ p(0)=0}
  \;\sum_{X : X_0 = c}\;\prod_{v \ne 0} [X_{p v} \incomp X_v]\,\|z(X_v)\| .
\]
Step C2 (\lc{pinnedClusterWeight_le_treeSumRaw}) dominates
$\mathrm{pinnedClusterWeight}(P,c,n)$ by
$\frac{1}{(n+1)!}\|z(c)\|\,\mathrm{treeSumRaw}(c,n,n)$: expand $|\Ur|$ over
spanning trees by Theorem~\ref{thm:penrose}, and embed each tree injectively
into admissible structures via its BFS parent and level maps (levels fit into
$\mathrm{Fin}(n+1)$ by \lc{connected_dist_lt_card}).  The master shell
recursion is then a statement about truncated Borel sums
$B_D^N(c) = \sum_{n \le N} \mathrm{treeSumRaw}(c,D,n)/n!$
(\lc{treeSumB}):
\[
  B_{D+1}^N(c) \;\le\; \exp\Bigl(\sum_{c'} [c \incomp c']\,\|z(c')\|\,B_D^N(c')\Bigr)
  \qquad (\lcmath{treeSumB_succ_le}),
\]
proved by cutting each admissible structure at the first shell (the children
of the root), factorizing the assignment sum over the subtrees hanging off
the shell (\lc{inner_factorization}, \lc{master_partition}), pricing the
enumeration of shells by multinomial coefficients
(\lc{card_blockData_mul_le}, \lc{per_k_bound}, \lc{rho_sum_le_price}), and
recognizing the exponential series in the $1/k!$-weighted size-vector sums
(\lc{treeSumRaw_succ_le}).  Induction on $D$ gives
$B_D^N(c) \le \Phi_D(c) \le e^{a(c)}$ (\lc{treeSumB_le_kpMajorant},
\lc{treeSumB_le_exp}); depth monotonicity of \lc{treeSumRaw} and the
factorial comparison $1/(n+1)! \le 1/n!$ assemble Theorem~\ref{thm:sharp}.

Why admissible parent/level pairs rather than trees?  Because they are a
\emph{decidable, finite, product-shaped} index (a filter on
$(\mathrm{Fin}(n+1)\to\mathrm{Fin}(n+1)) \times (\mathrm{Fin}(n+1)\to\mathrm{Fin}(D+1))$),
over-counting trees harmlessly (the inequality only needs an injection
from trees), while making ``the first shell'' and ``the subtree rooted at
$s$'' definable by simple arithmetic on levels (\lc{shellRoot},
\lc{shellFiber}, \lc{subtreeParent}, \lc{subtreeLev}).  Nothing in the
recursion manipulates a graph.

\section{The Mayer--Ursell inversion}\label{sec:mayer}

\begin{theorem}[Mayer--Ursell inversion]\label{thm:mayer}
\textup{[\lc{partition_eq_exp_clusterSum}, \file{KP/MayerInversion.lean}]}
Let $P$ be a finite polymer system whose cluster series is absolutely
convergent, i.e.\ $\sum_n \bigl\|\frac{1}{(n+1)!}\sum_X \Ur(X)\prod_i z(X_i)\bigr\|
< \infty$.  Then
\[
  \Xi(\Pol) \;=\; \exp\bigl(\mathrm{clusterSum}(P)\bigr).
\]
In particular (\lc{partition_eq_exp_clusterSum_of_kp}) this holds under the
criterion~\eqref{eq:kp}, and then $\Xi(\Pol) \ne 0$.
\end{theorem}
\begin{lstlisting}
theorem partition_eq_exp_clusterSum [Fintype P.Polymer]
    (hA : Summable fun n : ℕ => ‖(((n + 1).factorial : ℂ))⁻¹ *
      ∑ X : Fin (n + 1) → P.Polymer,
        (ursell P X : ℂ) * ∏ i, P.activity (X i)‖) :
    partition P (Finset.univ : Finset P.Polymer)
    = Complex.exp (clusterSum P)
\end{lstlisting}

This is the fundamental theorem of cluster expansions
(Friedli--Velenik Proposition~5.3, in the ``$\Xi = \exp$'' rather than
``$\log\Xi = $'' form, which avoids any branch of the complex logarithm).
Its formalization has a finite combinatorial heart and an analytic shell.

\paragraph{The partition identity.}
For a tuple $X$, call $X$ \emph{pairwise compatible}
(\lc{PairwiseCompatible}) if no two distinct coordinates are incompatible.
The indicator identity
$\sum_{E \subseteq E(G_X)} (-1)^{|E|} = [X\ \text{pairwise compatible}]$
(\lc{sum_neg_one_pow_eq_indicator}) is Mathlib's
\lc{Finset.sum_powerset_neg_one_pow_card}.  Grouping edge sets by the
partition of $\mathrm{Fin}\,n$ into connected components
(\lc{componentPartition}, a \lc{Finpartition} obtained from Mathlib's
reachability setoid) and factorizing each fiber over its blocks
(\lc{fiber_cp_factorization}) yields
\begin{equation}\label{eq:partition-identity}
  \sum_{\pi \,\in\, \mathrm{Finpartition}(\mathrm{Fin}\,n)}\;
  \prod_{B \in \pi} \Ur\bigl(X|_B\bigr) \;=\; [X\ \text{pairwise compatible}]
  \qquad (\lcmath{ursell_partition_identity}),
\end{equation}
where $X|_B$ is $X$ restricted to the block $B$ via Mathlib's order
isomorphism $\mathrm{Fin}\,|B| \simeq B$.  Summing~\eqref{eq:partition-identity}
against $\prod_j z(X_j)$ over all $N$-tuples, converting unordered partitions
into ordered ones with a $1/k!$ (\lc{sum_finpartition_eq_ordPartitions},
\lc{IsOrdPartition}) and counting the tuples with prescribed block sizes by
multinomials (\lc{card_ordPartition_mul}, \lc{sum_ordp_fiber_sizes}), one
arrives at the \emph{fully finite cluster form} of $\Xi$
(\lc{partition_univ_eq_cluster_layers}): $\Xi(\Pol)$ equals a finite sum over
total sizes $N \le \#\Pol$ of $1/k!$-weighted sums over size vectors
$(m_1,\dots,m_k)$, $\sum m_i = N$, $m_i \ge 1$, of products of normalized
per-block cluster weights $\frac{1}{m_i!}\sum_{Y:\mathrm{Fin}\,m_i\to\Pol}
\Ur(Y)\prod z(Y_l)$.  This is where the injectivity of pairwise compatible
tuples (\lc{PairwiseCompatible.injective}, from the hard core) turns the
tuple sum into the admissible-set sum $\Xi$.

\paragraph{The exponential of a \texttt{tsum}.}
Write $a_n$ for the $n$-th term of~\eqref{eq:clustersum}, so
$\mathrm{clusterSum} = \sum' a_n$.  Under absolute convergence,
\[
  \exp\Bigl(\sum'_n a_n\Bigr)
  \;=\; \sum'_{(k,f)\,\in\,\Sigma_k (\mathrm{Fin}\,k \to \mathbb{N})}
  \frac{1}{k!}\prod_{i} a_{f(i)}
  \qquad (\lcmath{exp_tsum_eq_tsum_H}),
\]
by Mathlib's power series for $\exp$ (\lc{NormedSpace.exp_eq_tsum_div}), the
identity $(\sum' a)^k = \sum'_{f}\prod_i a_{f(i)}$ (\lc{tsum_pow_eq_tsum_pi},
by induction on $k$ through \lc{tsum_mul_tsum_of_summable_norm}) and
\lc{Summable.tsum_sigma}.  The master family on the sigma type is absolutely
summable by comparison with the exponential series (\lc{summable_H}), and
regrouping it by total size $N = \sum_i (f(i)+1)$ (\lc{tsum_H_eq_tsum_layers})
produces exactly the finite layers of the cluster form, term by term
(\lc{layer_shift}).  Since the layers vanish for $N > \#\Pol$, the
\texttt{tsum} over $N$ is a finite sum and Theorem~\ref{thm:mayer} follows.
This is how the ``restricted'' inversion avoids infinite sums on the
partition-function side: $\Xi$ is a polynomial in the activities, the only
infinite object is $\mathrm{clusterSum}$, and absolute convergence is used
exactly once, to justify the sigma-type regrouping.

\paragraph{A consistency check.}
On the single-polymer system with $\|z\|<1$ the theorem specializes, via
$\Ur(K_{n+1}) = (-1)^n n!$ and the Mercator series
(\lc{mayer_log_series_complex}, from Mathlib's
\lc{Complex.hasSum_taylorSeries_log}), to $\mathrm{clusterSum} = \log(1+z)$
and $\Xi = 1 + z$ (\lc{clusterSum_singlePolymer_eq_log},
\lc{partition_singlePolymer_eq_exp}, \file{UrsellRecurrence.lean}).  This was
proved before the general theorem and served as the normalization check
between \lc{partition} and \lc{clusterSum}.

\paragraph{Zero-free polydisc (\texttt{ActivityDomain.lean}).}
The criterion reads $z$ only through $\|z\|$, so it is monotone under
pointwise domination (\lc{kpCriterion_withActivity_of_le}), and
Theorem~\ref{thm:mayer} gives:

\begin{corollary}
\textup{[\lc{partition_withActivity_ne_zero_of_kp}, \lc{partition_diskFamily_ne_zero_of_kp}]}
If $a$ satisfies~\eqref{eq:kp} for $P$ and $z'$ is any activity with
$\|z'(X)\| \le \|z(X)\|$ for all $X$, then $\Xi_{z'}(\Pol) \ne 0$.  In particular
$\Xi_{w z}(\Pol) \ne 0$ for all complex $w$ with $\|w\| \le 1$, and
$w \mapsto \Xi_{wz}(\Pol)$ is the explicit polynomial
$\sum_{S\ \mathrm{admissible}} w^{|S|}\prod_{X\in S} z(X)$
(\lc{partition_diskFamily_eq}).
\end{corollary}

\section{Restriction, tilting, tails, and animals}\label{sec:restrict}

The consumers need the engine in finite volume $\Lambda \subseteq \Pol$ and
need estimates whose constants do not see $\#\Pol$.  Four small modules
provide this.

\paragraph{Restriction (\texttt{Restriction.lean}).}
$P.\mathrm{restrict}\,\Lambda$ is the polymer system on the subtype
$\{X \,//\, X \in \Lambda\}$ with inherited incompatibility and activity.  Then
$\Xi_P(\Lambda) = \Xi_{P.\mathrm{restrict}\,\Lambda}(\mathrm{univ})$
(\lc{partition_restrict}), the criterion restricts with the restricted weight
(\lc{KPCriterion.restrict}, because the restricted criterion sums are
sub-sums), and hence \emph{every} finite-volume partition function is an
exponential: $\Xi_P(\Lambda) = \exp(\mathrm{clusterSum}(P.\mathrm{restrict}\,\Lambda))$
(\lc{partition_eq_exp_clusterSum_restrict}).  The ratio of two partition
functions is therefore $\exp$ of a \emph{difference} of cluster sums, and the
difference is exactly the tuple sum over clusters meeting $\Lambda^{c}$
(\lc{clusterSum_sub_restrict}), whose norm is bounded by the off-region
weight $\sum_n \mathrm{offRegionClusterWeight}(P,\Lambda,n)$
(\lc{norm_clusterSum_sub_restrict_le}).

\begin{theorem}[volume-free $Z$-ratio exponent]
\textup{[\lc{tsum_offRegionClusterWeight_le}]}
If $t > 0$ and $a$ satisfies the criterion for the $e^t$-scaled activities
$P.\mathrm{scaleActivity}(e^t)$, then
\[
  \sum_{n} \mathrm{offRegionClusterWeight}(P,\Lambda,n)
  \;\le\; t^{-1}\sum_{c \notin \Lambda} e^{t}\,\|z(c)\|\,e^{a(c)} .
\]
\end{theorem}
The proof pins each off-region cluster at a coordinate outside $\Lambda$
(\lc{offRegionClusterWeight_le_pinned}), pays the order factor $n+1$ by the
tilt $e^{t(n+1)} \ge t(n+1)$ (\lc{orderFactor_pinnedClusterWeight_le_tilt}),
and applies Theorem~\ref{thm:sharp} to the scaled system.  The algebra of the
cancellation, $\|\mathrm{Num}/Z\| \le \sum_i a_i e^{C_i}$ when
$Z = e^{\zeta}$ and $\|\zeta - \zeta_i\| \le C_i$ (\lc{norm_div_le_pinned_sum_exp}),
is stated separately so that the area-law gate can use it verbatim.

\paragraph{Tilting and size tails (\texttt{ClusterTail.lean}).}
$P.\mathrm{tilt}\,w$ rescales activities by $e^{w(c)}$ without touching
$\Ur$ (\lc{tilt_ursell} is \texttt{rfl}).  For a size function
$\mathrm{sz} : \Pol \to \mathbb{N}$ and the pinned weight restricted to
clusters of total size $\ge L$ (\lc{pinnedClusterWeightGE}):

\begin{theorem}[exponential size tail]
\textup{[\lc{kp_pinned_cluster_tail_bound}, \lc{pinned_cluster_tail_summable}]}
If $\varepsilon \ge 0$ and $a$ satisfies the criterion for
$P.\mathrm{tilt}(\varepsilon\,\mathrm{sz})$, then for all $c$, $L$, $N$
\[
  \sum_{n \le N} \mathrm{pinnedClusterWeightGE}(P,\mathrm{sz},c,L,n)
  \;\le\; e^{-\varepsilon L}\; e^{\varepsilon\,\mathrm{sz}(c)}\,\|z(c)\|\,e^{a(c)} .
\]
\end{theorem}
This is the decay mechanism behind the correlator and thermodynamic-limit
consumers: clusters connecting two distant regions are large.

\paragraph{Lattice animals, twice.}
The criterion~\eqref{eq:kp} for a lattice gas needs an entropy bound on
connected polymers.  Two independent theorems supply it, in different
generality.
\begin{itemize}
\item \lc{card_connectedPolymers_le} (\file{YangMills/L1_GibbsMeasure/ConnectedEntropy.lean}):
  on the periodic $d$-dimensional lattice of any size $N$, the number of
  connected plaquette sets of cardinality $n+1$ through a fixed plaquette is
  at most $(16d+1)^{2n}$ --- every such set is the range of a lazy closed
  walk of length $2n$ (\lc{exists_covering_lazyWalk}), and there are at most
  $(16d+1)^{2n}$ such walks because a plaquette touches at most $16d$ others
  (\lc{card_plaquettesTouching_le}).  The constant does not depend on $N$.
\item \lc{animal_card_le} (\file{YangMills/RG/AnimalTour.lean}): on any
  simple graph of maximum degree $\le \Delta$, any family of $S$-connected
  vertex sets of size $n$ containing a root has at most $\Delta^{2(n-1)}$
  members.  The proof peels the farthest vertex (\lc{exists_peel}: a
  max-distance vertex is never a cut vertex), builds a spanning closed walk of
  length $2(n-1)$ by splicing detours (\lc{exists_spanning_closed_walk},
  \lc{exists_detour_walk}), and counts walks by
  \lc{card_walks_length_le_degree_pow} (at most $\Delta^{\ell}$ walks of
  length $\ell$ from a vertex, from Mathlib's recursive
  \lc{finsetWalkLength}).  The cardinal form
  \lc{rooted_connected_card_le_pow} gives $c_n \le (\Delta^2)^n$ and feeds the
  geometric summability \lc{rooted_connected_weight_summable}.
\end{itemize}
\begin{lstlisting}
theorem animal_card_le [Fintype V] [DecidableRel G.Adj] (A : Finset (Finset V))
    {Δ n : ℕ} (hΔ : ∀ w, G.degree w ≤ Δ)
    (hA : ∀ S ∈ A, r ∈ S ∧ S.card = n ∧ ∀ x ∈ S, ∃ w : G.Walk r x, IsSWalk S w) :
    A.card ≤ Δ ^ (2 * (n - 1))
\end{lstlisting}

\section{How volume-uniformity is achieved, and the consumers}\label{sec:consumers}

The engine never leaves finite volume; volume-uniformity is a property of
the \emph{shape} of the bounds.  Concretely, for the connected lattice gas of
the strong-coupling expansion --- polymers are nonempty connected plaquette
sets, incompatibility is edge-support overlap, activities are gauge
expectations of $\prod_{p}(e^{-\beta\,\mathrm{pe}(\mathrm{hol}_p)}-1)$
(\lc{connectedLatticePolymerSystem}, \file{LatticePolymerSystem.lean}) ---
the chain is:

\begin{theorem}[volume-uniform criterion]
\textup{[\lc{connectedLatticePolymerSystem_kpCriterion_volumeUniform}, \file{ConnectedEntropy.lean}]}
Let $|\mathrm{pe}| \le B$, $t \ge 0$, $x = (e^{|\beta| B}-1)e^{t}$, and suppose
$(16d+1)^2 x < 1$ and $16d\cdot x/(1-(16d+1)^2 x) \le t$.  Then the weight
$a(c) = t\,|c|$ satisfies the criterion~\eqref{eq:kp} for the connected
lattice gas on the lattice of every linear size $N$.
\end{theorem}

The hypotheses mention $d$, $B$, $\beta$ and $t$ only.  The proof steers
every polymer $Y$ incompatible with $c$ through a plaquette $p \in Y$ touching
$c$ --- at most $16d$ per plaquette of $c$ (\lc{card_plaquettesTouching_le}),
hence at most $16d\,|c|$ in all --- and bounds each per-plaquette sum
$\sum_{Y \ni p} x^{|Y|}$ by the geometric series
$x/(1-(16d+1)^2x)$ using \lc{card_connectedPolymers_le}
(\lc{sum_connectedPolymers_through_le}).  Composed with
Theorem~\ref{thm:sharp} this gives absolute convergence of the cluster series
(\lc{connectedLatticeClusterSum_summable_volumeUniform}) and
$\|\mathrm{clusterSum}\| \le \sum_c \|z(c)\|e^{t|c|}$
(\lc{connectedLatticeClusterSum_norm_le_volumeUniform}) with no hypothesis on
$N$; the earlier uniform-majorant route needed $e\,t\,\#\{\text{plaquettes}\} < 1$
(\lc{connectedLatticeClusterSum_summable}), which is visible in the code as
the hypothesis the sharp version deleted.  The tilted criterion
(\lc{connectedLatticePolymerSystem_tilt_kpCriterion_volumeUniform}) and the
tail theorem give the volume-uniform size tail
\lc{connectedLattice_pinned_tail_volumeUniform}.

The three documented consumers are, briefly:
\begin{itemize}
\item \textbf{Area law} \cite{Eriksson2026AreaLaw}
  (\file{YangMills/L1_GibbsMeasure/RestrictedGate.lean}).  The region-restricted
  partition function of a weighted gas is written, by the truncation device
  $w \mapsto w\mathbf{1}_F$ and \lc{partition_eq_of_activity_eq_zero}, as a full
  partition function; \lc{partition_eq_exp_clusterSum_of_kp} and
  \lc{partition_eq_exp_clusterSum_restrict} turn $Z$ and each far factor into
  exponentials; \lc{norm_clusterSum_sub_restrict_le},
  \lc{tsum_offRegionClusterWeight_le} and \lc{norm_div_le_pinned_sum_exp}
  cancel $Z$ at a cost charged to the loop's perimeter; the pinned dichotomy
  is resummed by \lc{sum_pinned_dichotomy_le} using the area-extraction split
  \lc{sum_ite_pow_le}.  The headline
  \lc{normalized_exp_wilson_loop_area_law} is volume-uniform because every
  constant it inherits from the engine has the per-polymer shape above.
\item \textbf{Exponential clustering} (\file{TwoPlaquetteCorrelator.lean}).
  \lc{sun_two_plaquette_correlator_bound} bounds the normalized truncated
  two-plaquette correlation of $SU(N_c)$ lattice gauge theory by
  $e^{-\varepsilon k}$ for plaquettes at touching-graph distance $\ge 2k$,
  with constants in $d, N_c, \beta, s, t, \varepsilon$ only, through the
  size-tail theorem applied to the marked gas.
\item \textbf{Thermodynamic limit} \cite{Eriksson2026ThermoLimit}.  The Cauchy
  modulus of the finite-volume local expectations is the volume-uniform
  pinned tail (\file{ThermodynamicLimit.lean}); a genuine free-boundary
  exhaustion is compared with the periodic sequence through the same tails
  (\file{FreeBoundaryThermodynamicLimit.lean}); and the endpoint theorem
  \lc{tendsto_integerFreeBoundary}\hspace{0pt}\lc{ThermodynamicExpectation}
  (\file{IntegerThermodynamicLimit.lean}) converges the free-boundary
  exhaustion to the periodic infinite-volume state.
\end{itemize}
These papers document their own statements and windows; we list them only to
make precise which engine theorems they consume.

\section{The axiom oracle}\label{sec:oracle}

The repository's verification discipline is a per-theorem
\texttt{\#print axioms}.  Table~\ref{tab:oracle} lists every declaration
cited above that we treat as a headline, its file, and where its verbatim
oracle output is recorded.  Three record types occur: the verification
ledger \file{docs/VERIFICATION-LEDGER.md} (verbatim batch outputs at dated
checkpoints; the KP/Penrose release audit of 2026-06-10 at 8{,}209 jobs and
its Addenda~5--7 and~59 cover most of the table), the committed transcripts
\file{docs/oracle-transcripts/ORACLE-2026072*.txt} (about 4{,}600 lines each, one
entry per headline of \file{oracle_check.lean} at that date), and \file{oracle_check.lean}
itself.  Every recorded output is \oracle{}.

The revision of \file{oracle_check.lean} accompanying this paper (branch
\texttt{paper/kp-engine}) adds a section that prints the axioms of forty-two
further declarations from the table --- among them
\lc{abs_ursell_le_card_spanningTrees}, \lc{penrose_hfiber},
\lc{tree_walk_bound}, \lc{kp_pinned_cluster_bound}, \lc{kp_convergence_sharp},
\lc{ursell_partition_identity} and the volume-uniform lattice criterion ---
which were previously recorded only in the ledger.  For full disclosure: the
additions were made by static inspection of the sources (every name was
checked to exist at the cited file and line) and have \emph{not} yet been
re-run through the compiler on this machine; their oracle outputs in the
table are those recorded in the ledger at earlier checkpoints, and the
continuous-integration workflow does not compile Lean (it enforces the
zero-\texttt{sorry}, zero-axiom source invariants; see
Section~\ref{sec:repro}).  Re-running \texttt{lake env lean oracle\_check.lean}
at the current checkpoint is the one-line check that closes this gap.

\begin{footnotesize}
\begin{longtable}{@{}p{0.45\textwidth}p{0.30\textwidth}p{0.18\textwidth}@{}}
\caption{Headline declarations, files (under \texttt{YangMills/}; \texttt{L1} abbreviates \texttt{L1\_GibbsMeasure}), and oracle
records.  ``L'' = verification ledger (verbatim), ``T'' = committed oracle
transcripts of 2026-07-13/27, ``O'' = present in \texttt{oracle\_check.lean}
before this revision, ``O$^{+}$'' = added on branch \texttt{paper/kp-engine},
``R'' = batch statement in \texttt{REPRODUCIBILITY.md} (merged oracle script run
to exit~0 on \texttt{main} checkpoint \texttt{f0720ba7}; no per-declaration
line quoted).  Every recorded output is \oracle{}.}\label{tab:oracle}\\
\toprule
Declaration (\texttt{YangMills.} prefix omitted) & File & Record\\
\midrule
\endfirsthead
\toprule
Declaration & File & Record\\
\midrule
\endhead
\multicolumn{3}{l}{\emph{Substrate}}\\
\lc{KP.partition_union} & \file{KP/Basic.lean} & O$^{+}$\\
\lc{KP.ursell_eq_zero_of_not_isCluster} & \file{KP/Ursell.lean} & O$^{+}$\\
\lc{KP.ursellComplete_recurrence} & \file{KP/UrsellRecurrence.lean} & L (2026-06-10), O$^{+}$\\
\lc{KP.ursellComplete_eq} & \file{KP/UrsellRecurrence.lean} & L, O$^{+}$\\
\lc{KP.partition_singlePolymer_eq_exp} & \file{KP/UrsellRecurrence.lean} & L, O$^{+}$\\
\lc{KP.kp_activity_le}, \lc{KP.kp_neighbor_sum_le} & \file{KP/Criterion.lean} & O$^{+}$\\
\midrule
\multicolumn{3}{l}{\emph{Penrose tree-graph inequality}}\\
\lc{KP.interval_signed_sum} & \file{KP/PenroseScheme.lean} & L, O$^{+}$\\
\lc{KP.abs_signedSum_le_of_scheme} & \file{KP/PenroseScheme.lean} & L, O$^{+}$\\
\lc{KP.penroseTree_mem_spanningTrees} & \file{KP/PenroseBFS.lean} & L, O$^{+}$\\
\lc{KP.penrose_hfiber} & \file{KP/PenroseFiber.lean} & L, O$^{+}$\\
\lc{KP.abs_ursell_le_card_spanningTrees} & \file{KP/PenroseFiber.lean} & L, O$^{+}$\\
\lc{KP.abs_ursell_le_treeCount}, \lc{KP.treeCount_le_pow} & \file{KP/PenroseFiber.lean} & L, O$^{+}$\\
\lc{KP.succ_pow_le_exp_mul_factorial} & \file{KP/PenrosePrep.lean} & L, O$^{+}$\\
\midrule
\multicolumn{3}{l}{\emph{Tree-walk bounds and uniform KP}}\\
\lc{KP.tree_walk_bound} & \file{KP/WalkBound.lean} & L, O$^{+}$\\
\lc{KP.tree_walk_bound_vertexwise} & \file{KP/WalkBound.lean} & T, O\\
\lc{KP.tree_walk_bound_pinned} & \file{KP/PinnedWalk.lean} & O$^{+}$\\
\lc{KP.kp_per_size_bound}, \lc{KP.kp_convergence}, \lc{KP.kp_norm_clusterSum_le} & \file{KP/KPBound.lean} & L, O$^{+}$\\
\midrule
\multicolumn{3}{l}{\emph{Sharp KP}}\\
\lc{KP.clusterWeight_eq_sum_pinned} & \file{KP/PinnedCluster.lean} & O$^{+}$\\
\lc{KP.kpMajorant_le_exp} & \file{KP/SharpMajorant.lean} & O$^{+}$\\
\lc{KP.pinnedClusterWeight_le_treeSumRaw} & \file{KP/SharpShell.lean} & O$^{+}$\\
\lc{KP.treeSumB_succ_le}, \lc{KP.treeSumB_le_kpMajorant} & \file{KP/SharpShell.lean} & L (Add.~6), O$^{+}$\\
\lc{KP.kp_pinned_cluster_bound} & \file{KP/SharpShell.lean} & L (Add.~6), O$^{+}$\\
\lc{KP.pinned_cluster_summable_sharp} & \file{KP/SharpShell.lean} & L (Add.~6), O$^{+}$\\
\lc{KP.kp_clusterWeight_summable_sharp}, \lc{KP.kp_convergence_sharp}, \lc{KP.kp_norm_clusterSum_le_sharp} & \file{KP/SharpKP.lean} & L (Add.~6), O$^{+}$\\
\midrule
\multicolumn{3}{l}{\emph{Mayer--Ursell inversion and activity domain}}\\
\lc{KP.ursell_partition_identity} & \file{KP/MayerInversion.lean} & L (Add.~7), O$^{+}$\\
\lc{KP.partition_eq_exp_clusterSum} & \file{KP/MayerInversion.lean} & L (Add.~7), T, O\\
\lc{KP.partition_eq_exp_clusterSum_of_kp} & \file{KP/MayerInversion.lean} & L (Add.~7), T, O\\
\lc{KP.kpCriterion_withActivity_of_le} & \file{KP/ActivityDomain.lean} & T, O\\
\lc{KP.partition_withActivity_ne_zero_of_kp} & \file{KP/ActivityDomain.lean} & T, O\\
\lc{KP.partition_diskFamily_ne_zero_of_kp}, \lc{KP.partition_diskFamily_eq} & \file{KP/ActivityDomain.lean} & T, O\\
\midrule
\multicolumn{3}{l}{\emph{Restriction, tails}}\\
\lc{KP.KPCriterion.restrict}, \lc{KP.partition_eq_exp_clusterSum_restrict} & \file{KP/Restriction.lean} & L (8{,}237 jobs), O$^{+}$\\
\lc{KP.norm_clusterSum_sub_restrict_le}, \lc{KP.tsum_offRegionClusterWeight_le} & \file{KP/Restriction.lean} & L (8{,}237 jobs), O$^{+}$\\
\lc{KP.orderFactor_pinnedClusterWeight_le_tilt}, \lc{KP.tsum_finset_pinnedClusterWeight_le} & \file{KP/Restriction.lean} & T, O\\
\lc{KP.kp_pinned_cluster_tail_bound}, \lc{KP.pinned_cluster_tail_summable} & \file{KP/ClusterTail.lean} & O$^{+}$\\
\midrule
\multicolumn{3}{l}{\emph{Lattice animals and the volume-uniform criterion}}\\
\lc{card_connectedPolymers_le} & \file{L1/ConnectedEntropy.lean} & L (Add.~5), O$^{+}$\\
\lc{connectedLatticePolymerSystem_kpCriterion_volumeUniform} & \file{L1/ConnectedEntropy.lean} & L (Add.~5), O$^{+}$\\
\lc{connectedLatticeClusterSum_summable_volumeUniform} & \file{L1/ConnectedEntropy.lean} & L (Add.~6), O$^{+}$\\
\lc{connectedLatticeClusterSum_norm_le_volumeUniform} & \file{L1/ConnectedEntropy.lean} & L (Add.~6), O$^{+}$\\
\lc{connectedLattice_pinned_tail_volumeUniform} & \file{L1/ConnectedEntropy.lean} & O$^{+}$\\
\lc{RG.card_walks_length_le_degree_pow} & \file{RG/AnimalCount.lean} & T, O\\
\lc{RG.animal_card_le} & \file{RG/AnimalTour.lean} & L (Add.~59), T, O\\
\lc{RG.rooted_connected_card_le_pow} & \file{RG/AnimalTour.lean} & T, O\\
\midrule
\multicolumn{3}{l}{\emph{Consumers (documented elsewhere)}}\\
\lc{normalized_exp_wilson_loop_area_law} & \file{L1/RestrictedGate.lean} & L, T, O\\
\lc{sun_two_plaquette_correlator_bound} & \file{L1/TwoPlaquetteCorrelator.lean} & T, O\\
\lc{tendsto_integerFreeBoundary}\allowbreak\lc{ThermodynamicExpectation} & \file{L1/IntegerThermodynamicLimit.lean} & R, O\\
\bottomrule
\end{longtable}
\end{footnotesize}

For the declarations marked only O$^{+}$ (eleven of them:
\lc{partition_union}, \lc{ursell_eq_zero_of_not_isCluster},
\lc{kp_activity_le}, \lc{kp_neighbor_sum_le}, \lc{tree_walk_bound_pinned},
\lc{clusterWeight_eq_sum_pinned}, \lc{kpMajorant_le_exp},
\lc{pinnedClusterWeight_le_treeSumRaw}, \lc{kp_pinned_cluster_tail_bound},
\lc{pinned_cluster_tail_summable} and
\lc{connectedLattice_pinned_tail_volumeUniform}) we have no verbatim
per-declaration transcript; what is recorded for them is that they compile
inside the sorry-free, axiom-free \texttt{YangMillsCore} build (latest recorded
green build: 8{,}463 jobs), and they are consumed by declarations whose oracle
\emph{is} recorded (a declaration depending on an extra axiom would propagate
it to every consumer).  We state this rather than assert a transcript we do
not have.

\section{What is not formalized}\label{sec:not}

In the author's usual convention, the following are explicitly \emph{not}
claimed.
\begin{itemize}
\item \textbf{Infinite polymer systems.}  Every theorem carries
  \texttt{[Fintype P.Polymer]}.  The infinite-volume statements of
  \cite{KoteckyPreiss1986,FriedliVelenik2017} for countable polymer sets are
  not formalized; volume-uniformity is achieved by the shape of finite-volume
  bounds and by the consumers' Cauchy arguments, not by an infinite-volume
  cluster expansion.
\item \textbf{The sharp criteria of Fern\'andez--Procacci and
  Bissacot--Fern\'andez--Procacci}~\cite{FernandezProcacci2007,BFP2010}, and
  the Dobrushin form of the criterion.  Only the Koteck\'y--Preiss form
  \eqref{eq:kp} is proved.
\item \textbf{Tree-graph identities} (Brydges--Kennedy, Brydges--Federbush,
  Abdesselam--Rivasseau) and the Penrose \emph{identity} (as opposed to
  inequality).  The formalized statement is the inequality
  $|\Ur| \le \#\Trees$, which suffices for convergence.
\item \textbf{Cayley's formula.}  The tree count used is
  $\mathrm{treeCount}(m+1) \le (m+1)^{m+1}$, not $(m+1)^{m-1}$; the exact
  child-factorial Catalan identity over spanning trees of the complete graph
  is a separate, documented result of the programme~\cite{Eriksson2026Catalan}
  and is not used by the engine.
\item \textbf{The correlation-function cluster expansion} (the cluster
  representation of truncated correlations, Ueltschi's Theorem~2).  The
  consumers obtain correlation decay through marked gases and the size-tail
  theorem instead.
\item \textbf{The $\log$ form.}  The inversion is stated as
  $\Xi = \exp(\mathrm{clusterSum})$; no statement selects a branch of
  $\log \Xi$.
\item \textbf{Analytic continuation in the fugacity.}  The zero-free polydisc
  is proved; that $w \mapsto \Xi(wz)$ is entire and that $\log \Xi$ is analytic
  on the polydisc are on the frontier of the satellite
  \texttt{lean-zero-free-regions}, not in this layer.
\item \textbf{Cluster expansions with holes / renormalization-group inputs.}
  The RG lane (\file{YangMills/RG/}) consumes the vertexwise walk bound and the
  animal count, but its Ba\l{}aban--Dimock source estimates are carried as
  explicit theorem hypotheses (\file{HYPOTHESIS_FRONTIER.md}) and are not
  claimed here.
\item \textbf{Nothing about the continuum.}  As every status document of the
  repository states, no continuum limit, no Osterwalder--Schrader
  reconstruction and no continuum mass gap is formalized or claimed.
\end{itemize}

\section{Reproduction}\label{sec:repro}

All artifacts are in the public repository
\begin{center}
\url{https://github.com/lluiseriksson/THE-ERIKSSON-PROGRAMME}
\end{center}
The engine is \file{YangMills/KP/} (35 modules); its lattice-side consumers
are under \file{YangMills/L1_GibbsMeasure/}; the graph-theoretic animal count
is \file{YangMills/RG/AnimalCount.lean} and \file{YangMills/RG/AnimalTour.lean}.
The toolchain is pinned: Lean \texttt{v4.29.0-rc6} (\file{lean-toolchain}) and
the Mathlib commit
\begin{center}\texttt{07642720480157414db592fa85b626dafb71355b}\end{center}
(\file{lakefile.lean} and \file{lake-manifest.json} agree; see
\file{REPRODUCIBILITY.md}).  To rebuild the verified core and print the
oracle:
\begin{lstlisting}
lake exe cache get                # prebuilt Mathlib oleans
lake build YangMillsCore          # latest recorded: 8463 jobs, green
lake env lean oracle_check.lean   # prints the axioms of every headline
python scripts/check_consistency.py   # zero sorry, zero core axioms
\end{lstlisting}
A full core build exceeds the free continuous-integration tier, so the
repository's CI workflow (\file{.github/workflows/ci.yml}) does \emph{not}
compile Lean; it runs the comment-aware scanner
\file{scripts/check_consistency.py}, which fails on any \texttt{sorry} in the
proof tree or any \texttt{axiom} declaration in the verified-core tree.  The
compiler builds and oracle runs are performed locally and recorded verbatim
in \file{docs/VERIFICATION-LEDGER.md} and \file{docs/oracle-transcripts/}.
This paper's sources are \file{papers/kp-engine/main.tex}; the branch
\texttt{paper/kp-engine} carries the paper, the \file{oracle_check.lean}
additions of Section~\ref{sec:oracle}, and (only on that branch) a
one-line change to the CI trigger so that the scanner also runs on
\texttt{paper/**} branches.  The document compiles with
\texttt{latexmk -pdf -interaction=nonstopmode main.tex}; the script
\file{check_listings_unicode.py} audits the Lean listings against the
\texttt{literate} table.

\section{Conclusion}\label{sec:conclusion}

The Penrose tree-graph inequality, the Koteck\'y--Preiss criterion in its
sharp per-polymer form, and the Mayer--Ursell inversion are now theorems of
Lean~4 over Mathlib, with the standard three-axiom oracle and no
\texttt{sorry}.  The formalization is faithful to the textbook mathematics
but not to its usual presentation: ordered tuples replace multisets,
coordinate pinning replaces membership pinning, admissible parent/level
pairs replace trees inside the shell recursion, and the exponential of the
cluster series is handled by a sigma-type regrouping rather than by
$\log$.  Each of these choices traded a conceptual step for a bookkeeping
step that Mathlib's \texttt{Finset} and \texttt{tsum} libraries could
discharge, and together they kept a 14{,}441-line development
oracle-clean at every checkpoint.  The layer is already load-bearing: the
volume-uniform area law and the thermodynamic limit of the programme are
its consumers.  Natural next steps are the infinite-polymer version of the
criterion, the Fern\'andez--Procacci sharpening, and extraction of the
abstract polymer model, the Ursell coefficient and Theorems~\ref{thm:penrose}
and~\ref{thm:mayer} into Mathlib-style generic contributions.

\subsection*{Acknowledgments}
The Lean formalization was developed interactively with Anthropic's Claude,
in accordance with the ACM authorship policy on AI-assisted work; every
theorem stated here was checked by the Lean~4 kernel, with the axiom oracle
outputs archived in the repository ledger and transcripts.

\begin{thebibliography}{99}

\bibitem{Penrose1963}
O.~Penrose,
\emph{Convergence of fugacity expansions for fluids and lattice gases},
J.~Math.\ Phys.\ \textbf{4} (1963), 1312--1320.

\bibitem{Penrose1967}
O.~Penrose,
\emph{Convergence of fugacity expansions for classical systems},
in: T.~A.~Bak (ed.), \emph{Statistical Mechanics: Foundations and
Applications}, Benjamin, New York, 1967, 101--109.

\bibitem{KoteckyPreiss1986}
R.~Koteck\'y and D.~Preiss,
\emph{Cluster expansion for abstract polymer models},
Comm.\ Math.\ Phys.\ \textbf{103} (1986), 491--498.

\bibitem{Brydges1986}
D.~C.~Brydges,
\emph{A short course on cluster expansions},
in: K.~Osterwalder and R.~Stora (eds.), \emph{Critical Phenomena, Random
Systems, Gauge Theories} (Les Houches 1984), North-Holland, 1986, 129--183.

\bibitem{Dobrushin1996}
R.~L.~Dobrushin,
\emph{Estimates of semi-invariants for the Ising model at low temperatures},
in: \emph{Topics in Statistical and Theoretical Physics}, Amer.\ Math.\ Soc.\
Transl.\ (2) \textbf{177} (1996), 59--81.

\bibitem{MiracleSole2000}
S.~Miracle-Sol\'e,
\emph{On the convergence of cluster expansions},
Physica A \textbf{279} (2000), 244--249.

\bibitem{Ueltschi2004}
D.~Ueltschi,
\emph{Cluster expansions and correlation functions},
Moscow Math.\ J.\ \textbf{4} (2004), 511--522.

\bibitem{Scott2005}
A.~D.~Scott and A.~D.~Sokal,
\emph{The repulsive lattice gas, the independent-set polynomial, and the
Lov\'asz local lemma},
J.~Stat.\ Phys.\ \textbf{118} (2005), 1151--1261.

\bibitem{ScottSokal2005}
A.~D.~Scott and A.~D.~Sokal,
\emph{On dependency graphs and the lattice gas},
Combin.\ Probab.\ Comput.\ \textbf{15} (2006), 253--279.

\bibitem{FernandezProcacci2007}
R.~Fern\'andez and A.~Procacci,
\emph{Cluster expansion for abstract polymer models.\ New bounds from an old
approach},
Comm.\ Math.\ Phys.\ \textbf{274} (2007), 123--140.

\bibitem{BFP2010}
R.~Bissacot, R.~Fern\'andez and A.~Procacci,
\emph{On the convergence of cluster expansions for polymer gases},
J.~Stat.\ Phys.\ \textbf{139} (2010), 598--617.

\bibitem{Faris2010}
W.~G.~Faris,
\emph{Combinatorics and cluster expansions},
Probab.\ Surv.\ \textbf{7} (2010), 157--206.

\bibitem{FriedliVelenik2017}
S.~Friedli and Y.~Velenik,
\emph{Statistical Mechanics of Lattice Systems: A Concrete Mathematical
Introduction},
Cambridge University Press, 2017 (Chapter~5).

\bibitem{Edmonds2023}
C.~Edmonds,
\emph{Formal probabilistic methods for combinatorial structures using the
Lov\'asz local lemma},
in: CPP 2024, ACM, 2024; arXiv:2310.00513.

\bibitem{Lean4}
L.~de~Moura and S.~Ullrich,
\emph{The Lean 4 theorem prover and programming language},
in: Automated Deduction -- CADE 28, LNCS 12699, Springer, 2021, 625--635.

\bibitem{mathlib}
The mathlib Community,
\emph{The Lean mathematical library},
in: CPP 2020, ACM, 2020, 367--381.

\bibitem{Eriksson2026Master}
L.~Eriksson,
\emph{The Master Map: An Audit-First, Attack-Resistant Navigation Guide to
the Unconditional Solution of the 4D $SU(N)$ Yang--Mills Existence and Mass
Gap Problem},
ai.viXra:2602.0096, 2026.

\bibitem{Eriksson2026Audit}
L.~Eriksson,
\emph{Mechanical Audit Experiments and Reproducibility Appendix for a
Companion-Paper Programme on 4D $SU(N)$ Yang--Mills Existence and Mass Gap},
ai.viXra:2602.0117, 2026.

\bibitem{Eriksson2026Catalan}
L.~Eriksson,
\emph{A Machine-Verified Bijective Proof of the Rooted Child-Factorial
Catalan Identity over Spanning Trees of the Complete Graph},
ai.viXra:2607.0001, 2026.

\bibitem{Eriksson2026AreaLaw}
L.~Eriksson,
\emph{A Machine-Checked Volume-Uniform Wilson-Loop Area Law via a Formalized
Cluster Expansion},
ai.viXra:2607.0005, 2026; repository \file{paper/area-law/}.

\bibitem{Eriksson2026ThermoLimit}
L.~Eriksson,
\emph{A Machine-Checked Thermodynamic Limit for Local Lattice Gauge Gibbs
States},
ai.viXra:2607.0091, 2026; repository \file{papers/thermodynamic-limit-kp/}.

\bibitem{Eriksson2026C1}
L.~Eriksson,
\emph{Rooted-tree majorants for cluster expansions with holes} (C1 paper),
repository \file{papers/c1-rooted-tree-majorants/}, 2026.

\end{thebibliography}

\end{document}
